% ============================================================
% Chapter 1 -- Revised Introduction (Deep Literature Version)
% Oxford Academic / IMS Journal Style
% amsart + mathptmx + natbib + booktabs + lineno
% ============================================================

\documentclass[reqno, 12pt]{amsart}

\usepackage{mathptmx}
\usepackage[T1]{fontenc}
\usepackage[utf8]{inputenc}
\usepackage[a4paper, top=2.5cm, bottom=2.5cm, left=2.8cm, right=2.8cm]{geometry}
\usepackage{setspace}
\doublespacing
\usepackage[left]{lineno}
\usepackage{amsmath, amssymb, amsthm}
\usepackage{mathtools}
\usepackage[round, colon, authoryear, sort]{natbib}
\usepackage{booktabs}
\usepackage{microtype}
\usepackage[dvipsnames]{xcolor}
\usepackage[colorlinks=true, linkcolor=NavyBlue,
            citecolor=NavyBlue, urlcolor=NavyBlue]{hyperref}

% --- Custom commands ---
\newcommand{\tauij}{\tau_{ij}}
\newcommand{\ZHY}{\textsc{zhy}}
\newcommand{\HY}{\textsc{hy}}
\newcommand{\doi}[1]{\href{https://doi.org/#1}{\texttt{doi:#1}}}

% --- Theorem environments ---
\theoremstyle{plain}
\newtheorem{theorem}{Theorem}[section]
\newtheorem{lemma}[theorem]{Lemma}
\newtheorem{proposition}[theorem]{Proposition}
\newtheorem{corollary}[theorem]{Corollary}
\theoremstyle{definition}
\newtheorem{definition}[theorem]{Definition}
\theoremstyle{remark}
\newtheorem{remark}[theorem]{Remark}

% --- Document metadata ---
\title[Asynchronous Covariance Estimation and AI Reliability]{%
  Correlation Methods in the Statistical Analysis of Financial Trading Data\\[0.3em]
  \normalsize Chapter 1 --- Introduction (Revised and Extended, April 2026)}

\author{Yang Wu Azzollini}

\address{Lady Margaret Hall, University of Oxford,
  Oxford OX2~6QA, United Kingdom}

\email{yangwu.azzollini@gmail.com}

\date{April 2026 \quad (Original DPhil thesis: Hilary Term, 2016)}

\subjclass[2020]{Primary 62M10; Secondary 62P20, 62F10, 68T07}

\keywords{Asynchronous covariance estimation, Hayashi--Yoshida estimator,
  Epps effect, high-frequency data, AI reliability monitoring,
  uncertainty quantification, lead-lag estimation,
  non-synchronous trading, microstructure noise}

\begin{document}

\begin{abstract}
This chapter introduces the central statistical problem addressed in
this thesis: the estimation of covariance between processes observed
at irregular, non-synchronous times, with explicit quantification of
uncertainty as a function of the observed overlap structure. The
original thesis, completed at the University of Oxford in 2016 under
the supervision of Professor Brian~D.\ Ripley and Dr Peter Clifford,
derived the first exact finite-sample closed-form variance expression
for the general class of weighted asynchronous covariance estimators,
and proved the optimality of a specific weighting scheme via a Lagrange
multiplier argument. This revised introduction situates those results
within the decade of subsequent literature---extending from
endogeneity corrections and high-dimensional inconsistency to intrinsic
bias analysis---and introduces a new application domain not considered
in the original work: the reliability monitoring of artificial
intelligence systems in high-stakes settings. The mathematical
structure of the AI reliability problem is shown to be identical to
that of asynchronous covariance estimation in high-frequency finance,
providing a novel bridge between two currently disjoint literatures
with concrete implications for medical AI deployment.
\end{abstract}

\maketitle
\linenumbers

\tableofcontents

\bigskip\noindent\textit{Note on revision.}
This introduction has been substantially rewritten from the 2016 DPhil
version to reflect developments in the asynchronous covariance
estimation literature (2016--2026) and the emergence of AI reliability
monitoring as a major application domain. The mathematical results of
Chapters~4--6 are unchanged from the original thesis. All contribution
claims have been carefully verified against the current literature.

% ============================================================
\section{The Statistical Challenge of Asynchronous\\
High-Frequency Data}
\label{sec:challenge}
% ============================================================

\subsection{The electronic trading landscape}

Financial markets have undergone a fundamental transformation in the
decade since this thesis was first written. The migration from floor
trading to fully electronic execution, already well advanced by 2016,
is now essentially complete across all major asset classes. Trade and
quote data arrive at microsecond resolution across hundreds of
instruments simultaneously, each stream governed by its own stochastic
arrival process, each contaminated by microstructure noise of varying
character. The fundamental statistical challenge---how to extract
reliable estimates of correlation from these noisy, irregularly spaced,
asynchronous streams in real time---has not been resolved by the passage
of time. If anything, it has become more acute.

\subsection{The Epps effect and its origins}

The central difficulty is one of synchronisation. Two instruments, $X$
and $Y$, are observed at transaction times determined by independent
point processes. The observation times of $X$ and $Y$ do not coincide:
there is no natural common clock. When one attempts to estimate
covariance using returns computed over regular calendar-time intervals,
a severe downward bias emerges that intensifies as the sampling
frequency increases. This phenomenon was first documented empirically
by \citet{epps1979}, who showed using intraday transaction data from
the New York Stock Exchange that pairwise stock return correlations
decrease systematically as the sampling interval shortens---a finding
that has since been confirmed across equities, foreign exchange, and
futures markets \citep{toth2009, reno2003}.

The origins of the Epps effect have been studied extensively. Two
primary causes have been identified: the non-synchronicity of
transaction times across assets \citep{reno2003}, and the existence
of lead-lag relationships between returns, which cause correlations
computed over short intervals to underestimate the true underlying
comovement \citep{dejong1997}. A third, more subtle factor---the
finite reaction time of market participants to new information---has
been shown to impart a characteristic time scale to the phenomenon
that does not diminish even as trading frequency increases tenfold
\citep{toth2009}. The practical implication is unambiguous: naive
calendar-time synchronisation of asynchronous data introduces
systematic bias that cannot be eliminated by simply sampling more
frequently.

\subsection{The pre-synchronous era and its limitations}

The earliest approach to the non-synchronous trading problem was that
of \citet{scholes1977}, who noted that ordinary least squares estimates
of the market model are biased and inconsistent when stocks trade at
different times, and proposed lead-lag adjustments to the realised
covariance estimator. This line of work---incorporating lagged
cross-returns into the estimator---mitigates but does not eliminate
the bias, since the optimal number of lags is unknown and the
corrections introduce additional noise. The alternative approach of
imposing a regular grid by previous-tick interpolation or linear
interpolation before computing returns introduces its own biases,
particularly for illiquid assets with infrequent trades
\citep{barndorff2011}.

\subsection{The Hayashi--Yoshida estimator}

The estimator of \citet{hayashi2005} addressed this problem at its
root by operating directly on the raw asynchronous data, using the
overlap structure of successive price increments rather than requiring
any synchronisation of the original observations. If $\Delta P^{(1)}_i$
denotes the $i$-th increment of the first price process and
$\Delta P^{(2)}_j$ the $j$-th increment of the second, the estimator
sums all products $\Delta P^{(1)}_i \cdot \Delta P^{(2)}_j$ for which
the corresponding time intervals overlap:
\begin{equation}
  \widehat{[P^{(1)},P^{(2)}]}_{\HY}
  = \sum_{i,j} \Delta P^{(1)}_i \cdot \Delta P^{(2)}_j
    \cdot \mathbf{1}\!\left\{I_i \cap J_j \neq \emptyset\right\},
  \label{eq:HY}
\end{equation}
where $I_i = (\tau^{(1)}_{i-1}, \tau^{(1)}_i]$ and
$J_j = (\tau^{(2)}_{j-1}, \tau^{(2)}_j]$ are the observation intervals
for the two assets. This estimator is consistent for the true
integrated covariance as the mesh of observation times tends to zero
\citep{hayashi2005}, asymptotically normal \citep{hayashi2008}, and
free of the biases introduced by synchronisation \citep{griffin2011}.
It was a substantial advance.

\subsection{The critical gap: uncertainty quantification}

However, a critical question remained unanswered in the original
formulation:
\begin{quote}
\emph{What is the exact variance of this estimator, and how does it
depend on the observed overlap structure?}
\end{quote}
Without a tractable answer to this question, practitioners have no
principled basis for uncertainty quantification: they cannot know how
much to trust any particular correlation estimate given the specific
pattern of trades observed. \citet{griffin2011} derived closed-form
bias and variance expressions for the basic \HY{} estimator under the
specific assumption of Poisson sampling with independent i.i.d.\
microstructure noise---a result of considerable practical value, but
one that depends critically on the Poisson distributional assumption
for the observation times.

This thesis addresses the general case. In Chapter~5 we derive an
exact, finite-sample, closed-form expression for the variance of a
general class of weighted asynchronous covariance estimators, as an
explicit function of the observed overlap structure $\tauij$. This
result is \emph{not asymptotic} and requires \emph{no distributional
assumptions} on the arrival process of observations. We prove that a
specific set of weights---optimal in the sense of minimising this
variance within the class of linear estimators---can be derived
analytically via a Lagrange multiplier argument. We further show that
the resulting estimator admits a recursive update formula, processing
each new observation in $O(1)$ time, making it suitable for real-time
implementation.

% ============================================================
\section{A Decade of Subsequent Literature}
\label{sec:literature_fin}
% ============================================================

The literature on asynchronous covariance estimation has advanced
substantially since 2016, and intellectual honesty requires careful
positioning of the present work within it.

\subsection{Noise and asynchronicity: the bivariate case}

Several major papers have extended the theoretical foundations. The
multivariate realised kernel estimator of \citet{barndorff2011}
achieves consistency, positive semi-definiteness, and robustness to
measurement noise simultaneously---a combination that had not
previously been achieved---using a refresh-time synchronisation scheme
that discards a fraction of the data to impose a common observation
grid. \citet{aitsahalia2010} proposed a quasi-maximum likelihood
approach for the bivariate case that handles both noise and
asynchronicity without discarding observations, obtaining a consistent
and efficient estimator under Gaussian noise. These represent the main
competing approaches to the problem and form the natural benchmark
against which the results of this thesis should be assessed.

\subsection{Endogeneity of sampling times}

A different line of extension addresses the case where observation
times are endogenous---that is, where the timing of trades depends on
the price process itself, as would arise in models of informed trading
or strategic order submission. \citet{robert2012} established
volatility and covariation estimators under a general endogenous
microstructure model. \citet{koike2014} and \citet{koike2015}
extended the pre-averaged \HY{} estimator to this setting, deriving
an optimal weighting scheme for the pre-averaging step under time
endogeneity. These results are important context for Chapter~5:
our optimal weight derivation is complementary to, but distinct from,
the Koike weights, which address a different source of optimisation.

\subsection{High-dimensional inconsistency}

\citet{chakrabarti2022} established that the \HY{} estimator is
inconsistent in the high-dimensional regime---when the number of
assets $p$ grows proportionally with the number of observations
$n$---and characterised the limiting spectral distribution of the
resulting covariance matrix estimate. This result does not affect the
finite-sample variance formula of Chapter~5, which applies in the
bivariate case, but it is important context for understanding the
scope of applicability of the \HY{} framework.

\subsection{Intrinsic bias}

\citet{georgiadis2024} identified a previously overlooked intrinsic
bias in the \HY{} estimator arising from its telescoping structure:
certain observation points are effectively double-counted when
intervals overlap at their endpoints, causing the estimator to cancel
out relevant data points. This ``formulaic bias'' is present even in
the noise-free case and is distinct from the microstructure bias
studied by \citet{griffin2011}. The conditions under which it arises
and its magnitude under Poisson sampling are characterised. The
relationship between this bias and the weighted estimator of Chapter~5
is an open question.

\subsection{Position of this thesis}

Against this background, the contribution of Chapter~5 is positioned
as follows. The \citeauthor{griffin2011} variance formula is derived
under Poisson sampling---a specific distributional model. The formula
of Theorem~5.1 is \emph{non-asymptotic and model-free}: it holds for
any realised observation times $\tauij$, without distributional
assumptions on the arrival process. Furthermore, this thesis treats
the \emph{general class of weighted estimators} of the form
\begin{equation}
  \widehat{C}_w = \sum_{i,j} w_{ij}\,
    \Delta P^{(1)}_i \cdot \Delta P^{(2)}_j
    \cdot \mathbf{1}\!\left\{I_i \cap J_j \neq \emptyset\right\},
\end{equation}
and proves that the weights $w_{ij} = \tauij/(\tau_i^{(1)}\tau_j^{(2)})$,
where $\tauij = |I_i \cap J_j|$, $\tau_i^{(1)} = |I_i|$,
and $\tau_j^{(2)} = |J_j|$, minimise the exact finite-sample variance.
These two features---finite-sample exactness for arbitrary observation
times, and proved optimality of the weighting scheme---distinguish the
results of Chapter~5 from all existing work of which the author is
aware.

% ============================================================
\section{A New Context: Artificial Intelligence\\
Reliability Monitoring}
\label{sec:AI}
% ============================================================

\subsection{The mathematical parallel}

In the decade since this thesis was written, the statistical problems
it addresses have acquired a new and urgent application domain: the
reliability monitoring of artificial intelligence systems deployed in
high-stakes environments. This connection was not anticipated in the
original formulation, but it is not accidental. The mathematical
structure of the problem is identical.

Modern AI systems---large language models, medical diagnostic tools,
autonomous decision systems---produce outputs at irregular intervals
determined by query arrival times and computational latency. Multiple
AI components operating in parallel produce correlated output streams
that arrive asynchronously: a diagnostic model produces a prediction
at time $t_i$, a confidence estimator produces a score at time $t_j$,
a downstream verification system produces a check at time $t_k$. The
temporal structure of these outputs is precisely that of the
high-frequency financial data studied in this thesis: point process
observations, irregular spacing, correlation structure of fundamental
practical importance.

\subsection{The reliability gap in current AI monitoring}

The reliability question in AI has the same statistical character as
the correlation question in finance: \emph{when can you trust the
relationship between two streams, and how confident can you be in that
assessment given the specific observations available?} Current AI
monitoring practice offers no satisfactory answer. Production systems
rely on heuristic tools---Population Stability Index (PSI),
Kolmogorov--Smirnov tests, autoencoder reconstruction error---designed
for batch, synchronous, single-stream monitoring. They provide no
provable uncertainty bounds, cannot handle asynchronous observation
natively, and cannot detect the emergence of systematic lead-lag
relationships between correlated outputs.

The deficiencies of these approaches have been documented. Statistical
frameworks for assessing the reliability of generative AI have been
identified as a fundamental unsolved problem \citep{dobriban2025}.
Methods for detecting hallucinations in large language models through
semantic entropy \citep{farquhar2024} and for uncertainty
quantification in clinical AI \citep{lindenmeyer2025} represent recent
progress, but neither addresses the problem of real-time correlation
monitoring between asynchronous output streams with provable
uncertainty bounds.

\subsection{Medical AI and the stakes of reliability}

The stakes are highest in medicine. The work of \citet{singhal2023}
on Med-PaLM demonstrated that a large language model could exceed a
passing score on the US Medical Licensing Examination for the first
time; \citet{singhal2025} showed Med-PaLM~2 achieving performance
comparable to physician-level answers on nine of nine clinical quality
axes. These results establish the \emph{capability} of AI in clinical
settings. What they do not establish---and what no current framework
provides---is a principled method for monitoring whether those
capabilities are being reliably exercised in live deployment.

A radiologist using an AI diagnostic assistant, a clinician relying
on an AI-augmented monitoring system in an intensive care unit, a
risk manager using an AI-driven early warning system: none of them
currently have access to mathematically principled answers to the
question of how confident they should be in a given output, given what
the system has been doing recently. The gap between AI
\emph{capability} benchmarks and AI \emph{reliability} monitoring
tools has been identified as a critical obstacle to safe deployment
\citep{lindenmeyer2025}.

\subsection{The framework of this thesis as a solution}

The framework developed in this thesis provides the mathematical
foundation for addressing this gap. The $\tauij$ overlap structure
that determines the precision of a financial covariance estimate has
an exact analogue in the temporal relationship between AI output
streams. The recursive update formula that enables real-time financial
monitoring is equally applicable to real-time AI reliability
monitoring. The lead-lag estimation methodology of Chapter~6 provides
a principled tool for detecting systematic delays between correlated
AI outputs---a failure mode in which one component of an AI pipeline
begins to lag another in a way that signals emergent unreliability.
This failure mode has no current detection method in the literature.

% ============================================================
\section{Structure of This Work}
\label{sec:structure}
% ============================================================

Chapter~2 provides the financial background: the sources and
characteristics of high-frequency data, the electronic trading
environment, and the role of latency in constraining the class of
feasible statistical estimators. The key insight is that computational
speed places hard constraints on what is statistically possible
in real time---a constraint that is equally relevant in AI monitoring
settings.

Chapter~3 introduces the foundational computational tools:
exponentially weighted moving averages, dynamic linear models,
Kalman filtering. The emphasis throughout is on recursive formulations
that can be updated in $O(1)$ time per new observation, rather than
batch methods requiring $O(n)$ recomputation.

Chapter~4 addresses univariate volatility estimation for irregularly
spaced noisy data, building on the subsampling framework of
\citet{zhou1996}. We derive explicit expressions for the variance of
a class of estimators and establish infill consistency---a property
that will be mirrored in the bivariate setting of Chapter~5.

Chapter~5---the central theoretical contribution---develops the
extended \ZHY{} covariance estimator. The main results are:
\begin{enumerate}
  \item Theorem~5.1: exact finite-sample closed-form variance
    expression as an explicit function of $\tauij$;
  \item optimality proof: the weights
    $w_{ij} = \tauij/(\tau_i^{(1)}\tau_j^{(2)})$ minimise this
    variance in the class of linear estimators (Lagrange multiplier
    argument);
  \item recursive update formula: $O(1)$ update per new observation.
\end{enumerate}

Chapter~6 develops the maximum likelihood estimator for lead-lag delay
in the asynchronous setting, extending the contrast-optimisation
approach of \citet{hoffmann2013}. We demonstrate its application to
high-frequency financial instruments and discuss its implications for
AI reliability monitoring.

The conclusion identifies open problems---in particular, the
development of a sequential change detection framework for asynchronous
cross-stream correlation with finite-sample Type~I and Type~II error
guarantees---and outlines the research programme that follows from
this work.

% ============================================================
\section{The Broader Significance}
\label{sec:significance}
% ============================================================

The question this thesis addresses---how to know when to trust a
correlation estimate computed from asynchronous noisy data---has
acquired a significance in 2026 that could not have been fully
anticipated in 2016. AI systems are now deployed in medical diagnosis,
drug discovery, and financial risk management, domains where errors
have direct human consequences. The deployment has outpaced the
development of rigorous tools for monitoring whether those systems are
behaving as expected.

The dominant narrative in AI research concerns capability: making
systems more powerful, more accurate, more general. The question of
reliability---\emph{when can you trust the system enough to act on its
output}---receives far less rigorous attention. Current monitoring
practice relies on heuristic thresholds whose statistical properties
are either unknown or depend on distributional assumptions that
production data routinely violates.

The mathematical framework of this thesis does not solve that problem
in its full generality. But it provides a rigorous foundation for one
critical component of the answer: real-time monitoring of cross-stream
correlation with provable uncertainty bounds, for genuinely asynchronous
point process observations. That is a narrow contribution, but a
genuine one. In a field where the gap between what is claimed and what
is proved has grown dangerously wide, narrowness paired with rigour is
not a weakness. It is exactly what is needed.

% ============================================================
\begin{thebibliography}{99}

\bibitem[\protect\citeauthoryear{{A\"{i}t-Sahalia}, Fan and Xiu}%
  {{A\"{i}t-Sahalia} et~al.}{2010}]{aitsahalia2010}
\textsc{A\"{i}t-Sahalia, Y., Fan, J. and Xiu, D.} (2010).
High-frequency covariance estimates with noisy and asynchronous
financial data.
\textit{Journal of the American Statistical Association}
\textbf{105}(492), 1504--1517.
\doi{10.1198/jasa.2010.tm10163}.

\bibitem[\protect\citeauthoryear{Barndorff-Nielsen, Hansen, Lunde and
  Shephard}{Barndorff-Nielsen et~al.}{2011}]{barndorff2011}
\textsc{Barndorff-Nielsen, O.~E., Hansen, P.~R., Lunde, A.
  and Shephard, N.} (2011).
Multivariate realised kernels: Consistent positive semi-definite
estimators of the covariation of equity prices with noise and
non-synchronous trading.
\textit{Journal of Econometrics} \textbf{162}(2), 149--169.
\doi{10.1016/j.jeconom.2010.07.009}.

\bibitem[\protect\citeauthoryear{Chakrabarti and Sen}{Chakrabarti and
  Sen}{2022}]{chakrabarti2022}
\textsc{Chakrabarti, A. and Sen, R.} (2022).
Limiting spectral distribution of high-dimensional {Hayashi--Yoshida}
estimator of integrated covariance matrix.
\textit{arXiv preprint} arXiv:2201.00119.

\bibitem[\protect\citeauthoryear{De~Jong and Nijman}{De~Jong and
  Nijman}{1997}]{dejong1997}
\textsc{De~Jong, F. and Nijman, T.} (1997).
High frequency analysis of lead-lag relationships between financial
markets.
\textit{Journal of Empirical Finance} \textbf{4}(2--3), 259--277.

\bibitem[\protect\citeauthoryear{Dobriban}{Dobriban}{2025}]{dobriban2025}
\textsc{Dobriban, E.} (2025).
Statistical methods in generative {AI}.
\textit{arXiv preprint} arXiv:2509.07054.

\bibitem[\protect\citeauthoryear{Epps}{Epps}{1979}]{epps1979}
\textsc{Epps, T.~W.} (1979).
Comovements in stock prices in the very short run.
\textit{Journal of the American Statistical Association}
\textbf{74}(366), 291--298.

\bibitem[\protect\citeauthoryear{Farquhar, Kossen, Kuhn and Gal}%
  {Farquhar et~al.}{2024}]{farquhar2024}
\textsc{Farquhar, S., Kossen, J., Kuhn, L. and Gal, Y.} (2024).
Detecting hallucinations in large language models using semantic
entropy.
\textit{Nature} \textbf{630}(8017), 625--630.
\doi{10.1038/s41586-024-07421-0}.

\bibitem[\protect\citeauthoryear{Georgiadis}{Georgiadis}{2024}]{georgiadis2024}
\textsc{Georgiadis, E.} (2024).
A note on asynchronous challenges: Unveiling formulaic bias and data
loss in the {Hayashi--Yoshida} estimator.
\textit{arXiv preprint} arXiv:2404.18233.

\bibitem[\protect\citeauthoryear{Griffin and Oomen}{Griffin and
  Oomen}{2011}]{griffin2011}
\textsc{Griffin, J.~E. and Oomen, R.~C.~A.} (2011).
Covariance measurement in the presence of non-synchronous trading and
market microstructure noise.
\textit{Journal of Econometrics} \textbf{160}(1), 58--68.
\doi{10.1016/j.jeconom.2010.03.015}.

\bibitem[\protect\citeauthoryear{Hayashi and Yoshida}{Hayashi and
  Yoshida}{2005}]{hayashi2005}
\textsc{Hayashi, T. and Yoshida, N.} (2005).
On covariance estimation of non-synchronously observed diffusion
processes.
\textit{Bernoulli} \textbf{11}(2), 359--379.
\doi{10.3150/bj/1116340299}.

\bibitem[\protect\citeauthoryear{Hayashi and Yoshida}{Hayashi and
  Yoshida}{2008}]{hayashi2008}
\textsc{Hayashi, T. and Yoshida, N.} (2008).
Asymptotic normality of a covariance estimator for nonsynchronously
observed diffusion processes.
\textit{Annals of the Institute of Statistical Mathematics}
\textbf{60}(2), 367--406.

\bibitem[\protect\citeauthoryear{Hoffmann, Rosenbaum and
  Yoshida}{Hoffmann et~al.}{2013}]{hoffmann2013}
\textsc{Hoffmann, M., Rosenbaum, M. and Yoshida, N.} (2013).
Estimation of the lead-lag parameter from non-synchronous data.
\textit{Bernoulli} \textbf{19}(2), 426--461.
\doi{10.3150/11-BEJ407}.

\bibitem[\protect\citeauthoryear{Huang, Yang, Zhu and Chen}{Huang
  et~al.}{2024}]{huang2024triadic}
\textsc{Huang, Y., Yang, K., Zhu, Z. and Chen, L.} (2024).
{Triadic-OCD}: Asynchronous online change detection with provable
robustness, optimality, and convergence.
In \textit{Proceedings of the 41st International Conference on Machine
Learning}, Vol.~235 of \textit{Proceedings of Machine Learning
Research}, pp.~20382--20412. PMLR.

\bibitem[\protect\citeauthoryear{Koike}{Koike}{2014}]{koike2014}
\textsc{Koike, Y.} (2014).
Limit theorems for the pre-averaged {Hayashi--Yoshida} estimator with
random sampling.
\textit{Stochastic Processes and their Applications}
\textbf{124}(8), 2699--2753.

\bibitem[\protect\citeauthoryear{Koike}{Koike}{2015}]{koike2015}
\textsc{Koike, Y.} (2015).
Time endogeneity and an optimal weight for the pre-averaged
{Hayashi--Yoshida} estimator.
\textit{Statistical Inference for Stochastic Processes}
\textbf{18}(3), 234--281.

\bibitem[\protect\citeauthoryear{Lindenmeyer, Blattmann, Franke,
  Neumuth and Schneider}{Lindenmeyer et~al.}{2025}]{lindenmeyer2025}
\textsc{Lindenmeyer, A., Blattmann, M., Franke, S., Neumuth, T.
  and Schneider, D.} (2025).
Towards trustworthy {AI} in healthcare: Epistemic uncertainty
estimation for clinical decision support.
\textit{Journal of Personalized Medicine} \textbf{15}(2), 58.
\doi{10.3390/jpm15020058}.

\bibitem[\protect\citeauthoryear{Ren\`{o}}{Ren\`{o}}{2003}]{reno2003}
\textsc{Ren\`{o}, R.} (2003).
A closer look at the {Epps} effect.
\textit{International Journal of Theoretical and Applied Finance}
\textbf{6}(1), 87--102.

\bibitem[\protect\citeauthoryear{Robert and Rosenbaum}{Robert and
  Rosenbaum}{2012}]{robert2012}
\textsc{Robert, C.~Y. and Rosenbaum, M.} (2012).
Volatility and covariation estimation when microstructure noise and
trading times are endogenous.
\textit{Mathematical Finance} \textbf{22}(1), 133--164.
\doi{10.1111/j.1467-9965.2010.00454.x}.

\bibitem[\protect\citeauthoryear{Scholes and Williams}{Scholes and
  Williams}{1977}]{scholes1977}
\textsc{Scholes, M. and Williams, J.} (1977).
Estimating betas from nonsynchronous data.
\textit{Journal of Financial Economics} \textbf{5}(3), 309--327.
\doi{10.1016/0304-405X(77)90041-1}.

\bibitem[\protect\citeauthoryear{Singhal et~al.}{Singhal
  et~al.}{2023}]{singhal2023}
\textsc{Singhal, K., Azizi, S., Tu, T. et~al.} (2023).
Large language models encode clinical knowledge.
\textit{Nature} \textbf{620}, 172--180.
\doi{10.1038/s41586-023-06291-2}.

\bibitem[\protect\citeauthoryear{Singhal et~al.}{Singhal
  et~al.}{2025}]{singhal2025}
\textsc{Singhal, K., Tu, T., Gottweis, J. et~al.} (2025).
Toward expert-level medical question answering with large language
models.
\textit{Nature Medicine} \textbf{31}(3), 943--950.
\doi{10.1038/s41591-024-03423-7}.

\bibitem[\protect\citeauthoryear{T\'{o}th and Kert\'{e}sz}{T\'{o}th
  and Kert\'{e}sz}{2009}]{toth2009}
\textsc{T\'{o}th, B. and Kert\'{e}sz, J.} (2009).
The {Epps} effect revisited.
\textit{Quantitative Finance} \textbf{9}(7), 793--802.

\bibitem[\protect\citeauthoryear{Zhou}{Zhou}{1996}]{zhou1996}
\textsc{Zhou, B.} (1996).
High-frequency data and volatility in foreign-exchange rates.
\textit{Journal of Business \& Economic Statistics}
\textbf{14}(1), 45--52.
\doi{10.1080/07350015.1996.10524628}.

\end{thebibliography}

\end{document}
